\chapter{Discovery across independent, latent, temporal and experimental settings}
\label{exec16:s2}

\section{Choose the observation experiment before the discovery engine}

A bank may want a causal diagram for several quite different reasons. It may
need to identify an adjustment set for an offer, investigate whether a funding
relationship survives a policy change, or decide which experiment would resolve
an ambiguity that matters financially. These uses do not imply the same
learning problem. A discovery engine can be useful while leaving the effect
unresolved; it can also orient more edges under assumptions that the bank has
no reason to accept. The proposal therefore treats discovery as a conditional
inference about a declared observation experiment, followed by a separate
decision about what the inferred object permits.

Let $V=(V_1,\ldots,V_p)$ denote the measured variables, $G$ a causal graph and
$P$ their observational distribution. The observed data may be independent
customer records, selected records conditional on inclusion $S=1$, trajectories
$V_{it}$ indexed by customer $i$ and physical time $t$, or observations from
regimes $e$ with different intervention targets. The symbol $P$ must be expanded
accordingly. A pooled table need not represent independent observations from
one distribution. Independence across customers does not establish independence
across their months, and differences across regimes need not be interventions
on the variables named by an analyst.

A discovery request should first state its admissible model family
$\mathcal M$, the observed law or laws and the graph object it seeks. Under
causal sufficiency and the relevant acyclicity, Markov and faithfulness
conditions, an observational conditional-independence procedure can aim at a
Markov equivalence class represented by a CPDAG. With hidden or selection
variables, the appropriate ancestral-graph representation and its guarantees
change. Temporal and interventional methods add further restrictions concerning
lags, regimes, mechanisms and sampling. The pcalg, Tetrad and causal-learn
comparisons make these differences concrete at their actual interfaces.
\citep{exec13_pcalg,exec13_java_tetrad,exec13_pydisc_causallearn,exec13_latentr_colombo2012}

There are then three separate outputs to consider. The first is an estimated
structural object, such as a CPDAG, PAG, fitted structural model or candidate
DAG. The second is the collection of causal answers compatible with that
object and the declared assumptions. The third is a statistical statement
about the uncertainty of the complete procedure. An estimated graph is not
automatically a confidence set for graphs. A collection of effects computed
from compatible DAGs is not automatically a confidence region. The inspected
pcalg IDA documentation, for example, describes a multiset of possible total
effects under its model assumptions. Keeping its multiplicities can be useful
for checking the computation, but counting repeated values does not establish
probabilities on causal worlds. \citep{exec13_pcalg,exec13_pcalg_reuse_primary}

These distinctions give a practical order for the software comparison. First
admit an observation experiment and causal model family. Then compare methods
that target the same structural object. Inspect what the implemented method
actually returns, including unresolved marks, conflicts and failures. Finally
connect that result to a stated banking action and loss. The preceding
R, Java and Python comparisons supply the first implementation examples;
the later latent, temporal and interventional treatments extend that evidence.
The present chapter connects these comparisons through the observation law
and develops the consequences for the banking decision.

\section{Independent observations: three ways to obtain direction}

Constraint-based methods use conditional independences to remove adjacencies
and orient features shared by compatible graphs. Their empirical behavior
depends on the independence test, conditioning sets, missing-data treatment
and orientation procedure. The test must correspond to the observed variables
and sampling law. A Fisher-Z partial-correlation test for a Gaussian model,
for example, is not a generic nonlinear conditional-independence oracle.
The inspected Tetrad implementation exposes missing-data, rank and
degrees-of-freedom choices that affect its returned decisions. Its PC path
and its score-based FGES path therefore need separate admission profiles even
though they come from the same software family. \citep{exec13_java_tetrad}

Score-based search compares candidate parent sets or equivalence classes
through a specified model score. The relevant questions include the likelihood
family, score equivalence, penalty, search neighborhood, admissible constraints
and stopping behavior. A finite optimum of a numerical objective supplies no
general causal guarantee. A local maximum may be a useful candidate, but it
should remain a candidate with its objective and search settings attached.
The R and native comparisons distinguish equivalence-class search from a
selected representative and distinguish a package wrapper from the numerical
kernel it invokes. \citep{exec13_pcalg,exec13_cxx_xges,exec16_ct_pcalg}

Functional-model methods can extract direction from restrictions beyond
conditional independences. Linear non-Gaussian, additive-noise and differentiable
structural-model approaches should therefore be compared on their actual
assumptions, not merely on whether each returns a weighted matrix. A
non-Gaussian independence argument cannot be inherited by a Gaussian example
simply because the same optimizer accepts its data. Likewise, a small
acyclicity penalty or thresholded coefficient matrix is a numerical diagnostic;
the graph and the conditions behind an intervention interpretation still need
checking. The inspected LiNGAM, NOTEARS, DAGMA and Causica paths offer
different candidate computations and require different proposed acceptance
tests. \citep{exec13_pydisc_lingam,exec13_pydisc_notears,exec13_pydisc_dagma,exec13_pydisc_causica}

The proposal does not select one of these families universally. A sufficient
linear Gaussian model with credible independent units gives a transparent
starting point for constrained PC and score-search comparisons. A credible
non-Gaussian structural restriction may justify a separate functional-model
comparison. A concern about hidden risk changes the admissible graph class
before it changes the preferred implementation language. The analyst should
be able to inspect the reason for that change and the additional assumptions
needed to reverse it.

\subsection{Faithfulness and the business target can fail in different ways}

Consider an authored structural example with independent, centered,
unit-variance Gaussian disturbances $U,V,W$:
\[
 X=U,\qquad M=X+V,\qquad
 Y=(-1+\varepsilon)X+M+W.
\]
The variables are standardized indices; their labels do not assert a real
bank mechanism. Substitution gives $Y=\varepsilon X+V+W$. Thus
\[
 \operatorname{Cov}(X,Y)=\varepsilon,
 \qquad \frac{\partial E[Y\mid\operatorname{do}(X=x)]}{\partial x}
 =(-1+\varepsilon)+1=\varepsilon.
\]
At $\varepsilon=0$, $X$ and $Y$ are marginally independent even though the
structural graph contains a direct edge and a directed indirect path. A
discovery conclusion that interprets every independence as a graphical
separation would lose this cancellation structure. The direct effect is
$-1$, the indirect effect is $1$ and the total effect is zero. These are
different targets even though they arise from the same three equations.
\citep{exec16_s2_exact}

The example also shows why graph error and decision error should both be
reported. Suppose that changing $X$ by one unit has value USD~30 per unit of
$Y$ and costs USD~8. With $|\varepsilon|\leq 1/20$, its net value lies between
USD~$-9.5$ and USD~$-6.5$. Declining this action is robust within that
stipulated family, even while the mechanism decomposition is unresolved.
An intervention on $M$ would answer a different question and would require
its own operational interpretation and cost. More arrows are not necessarily
worth acquiring for the original action. \citep{exec16_s2_exact}

Near cancellation is also relevant to finite discovery. When an empirical
test has little power against a small nonzero association, treating a
nonrejection as a known absence can be consequential for other targets.
Changing the test threshold after observing the desired graph does not repair
that problem. The proposed result should retain the tests, separation sets,
search limits and unresolved features needed to explain the conclusion.
Stability under column order is one diagnostic, but it is not identification
or calibration. The order-independent literature explicitly separates stable
skeleton estimation, collider decisions and later orientation rules.
\citep{exec13_latentr_colombo2014,exec13_pydisc_causallearn}

\section{A common observational law can support different actions}

The following three-variable example makes the bridge from graph discovery to
decision explicit. All variables are centered. Within each model the listed
disturbances are independent Gaussian variables with the stated variances:
\begin{align*}
 F:\quad &X=U_X,\quad M=X+U_M,\quad Y=M+U_Y,\\
 R:\quad &Y=U_Y,\quad M=\tfrac23Y+U_M,\quad X=\tfrac12M+U_X,\\
 C:\quad &M=U_M,\quad X=\tfrac12M+U_X,\quad Y=M+U_Y.
\end{align*}
In the order $(U_X,U_M,U_Y)$ the variance vectors are respectively
$(1,1,1)$, $(1/2,2/3,3)$ and $(1/2,2,1)$.
Here $F$ is the forward chain, $R$ the reverse chain and $C$ the common-parent
model. The noise labels are local to each model. In the coordinate order
$(X,M,Y)$, all three observational distributions have covariance
\[
 \Sigma=\begin{pmatrix}1&1&1\\1&2&2\\1&2&3\end{pmatrix}.
\]
For example, in $C$, $\operatorname{Var}(X)=\tfrac14(2)+\tfrac12=1$,
$\operatorname{Cov}(X,Y)=\tfrac12(2)=1$ and
$\operatorname{Var}(Y)=2+1=3$. The other two factorizations give the same
matrix. As centered Gaussian laws they are therefore identical, rather than
merely having the same least-squares fit. Their conditional covariance of
$X$ and $Y$ given $M$ is $1-(1)(2)/2=0$.
\citep{exec16_s2_exact}

The three graphs have the same chain skeleton and no collider. If this
three-model family is the stipulated scientific scope, observations alone
leave all three possibilities. The effect of changing $X$ from zero to one
on $Y$ is one in $F$ and zero in both $R$ and $C$. The compatible effect set
is $\{0,1\}$; a convex interval $[0,1]$ is a conservative enclosure, not the
same exact set. Selecting one representative DAG would turn an unresolved
direction into a selected action effect. Conversely, a verified ordering
that truly rules out some directions is additional information and should
be recorded as such. Simultaneous quarterly summaries with convenient names
do not supply that ordering. \citep{exec16_s2_exact}

For a hypothetical offer, retain the stipulated USD~30 outcome scale and
USD~8 action cost. Net values are USD~22 in $F$ and USD~$-8$ in $R$ and $C$.
Different decision criteria can consequently recommend different actions.
The worst-case value of offering is $-8$, compared with zero for declining,
so a maximin rule declines. The largest regret from offering is $8$, while
the largest regret from declining is $22$, so a minimax-regret rule offers.
If one stipulates equal weights on the three models, expected net value is
$22/3-8/3-8/3=2$, and that expected-value rule also offers. The equal weights
are a scenario assumption, not a probability distribution learned by a CPDAG
procedure. A product should retain the distinction between the compatible
model set and the sponsor's decision criterion. \citep{exec16_s2_exact}

\section{Hidden variables and selection change the object being learned}

Unmeasured borrower risk, market shocks and selection into a reviewed sample
can invalidate a causally sufficient graph interpretation. Introducing a
hidden node in a diagram is not enough: the learning procedure and the
representation must preserve the relevant conditional-independence and
ancestral information. An ancestral graph can represent relationships after
marginalizing hidden variables. A partial ancestral graph records invariant
endpoint information across a class, with circles retaining unresolved marks.
Its directed edge is not generally a promise that an unconfounded regression
coefficient equals a direct causal effect. The inspected ancestral-graph
definitions and the software comparisons distinguish these meanings.
\citep{zhang2008,exec13_latentr_primary}

FCI and RFCI illustrate why even two outputs called PAGs need additional
method identity. The RFCI definition permits an adjacency with a weaker
interpretation than the corresponding FCI adjacency. Its graph can retain
edges that a more extensive search removes, and can correspond to more than
one Markov equivalence class. The oracle soundness statement assumes a
faithful distribution and the required conditional-independence information;
it is not a finite-sample confidence claim. The faster procedure is therefore
a substantive information/computation alternative, not a lossless substitute
whose edges can automatically be read under the stronger definition.
\citep{exec13_latentr_colombo2012}

Selection is equally concrete. Let $X$ and $U$ be independent fair binary
variables, let $Y=U$, and retain a record only when $S=1$, where
$S=\mathbf1\{X=U\}$. The four population configurations each have probability
$1/4$, while the retained sample consists of $(X,Y)=(0,0)$ and $(1,1)$ with
equal conditional probability. Its observed conditional contrast is one.
Yet replacing the equation for $X$ leaves $U$ and $Y$ unchanged, so the
population intervention contrast is zero. These are exact properties of a
generated law. They show why the inclusion event and target population belong
in the request. They do not claim that a latent-discovery routine can recover
the omitted population law from the selected table alone.
\citep{exec16_s2_exact}

The practical candidate portfolio therefore keeps pcalg FCI/RFCI, corresponding
Tetrad and causal-learn paths, and selected native or other-language latent
methods as distinct comparisons. Source-specific endpoint conversions,
conditioning bounds and failure semantics remain relevant. A PC conflict
represented by two arrowheads must not be silently relabeled as a latent
confounding conclusion. A latent mixture model or a supplied Bayesian network
with hidden states is also a different capability from learning an ancestral
equivalence class. The commercial and native latent comparisons develop
those alternatives without treating all hidden-node software as an FCI
replacement. \citep{exec13_pydisc_causallearn,exec13_latentr_primary,exec14_native_agrum,exec14_native_miic,exec13_latent_genie_structure,exec13_latent_mplus}

\section{Temporal discovery begins with a physical clock}

For a bank panel, write a candidate structural process as
\[
 V_{it}=B_0V_{it}+\sum_{\ell=1}^{L}B_\ell V_{i,t-\ell}+U_{it}.
\]
The ordering convention here is target by source: entry $(j,k)$ of $B_\ell$
describes the effect of coordinate $k$ at lag $\ell$ in the equation for
coordinate $j$. This is a proposed model, not a description of all temporal
packages. Its interpretation requires a physical time unit, a definition of
the contemporaneous interval, an admissible structure for $B_0$ and a law
for $U_{it}$. If $I-B_0$ is invertible, the reduced lag coefficient is
$(I-B_0)^{-1}B_\ell$. Predicting the reduced equation does not by itself
identify the structural contemporaneous matrix.

A customer boundary and a calendar gap matter before any conditional test.
Joining the final month of one customer to the first month of another creates
a fictitious transition. Treating a three-month gap as one ordinary lag
changes the physical model. A recorded variable can also be unavailable at
the action time even if its event time is earlier. Neither column order nor a
lagged dataframe constructor certifies these conditions. The R, Python, Java
and C++ temporal comparisons trace their different window, subject, regime
and state-label contracts. \citep{exec14_pt_contract,exec16_ct_pcalg,exec16_ct_agrum,exec16_ct_sdbn}

A small impulse calculation shows why the outcome horizon belongs beside the
graph. Consider a two-period response with a unit action at the first period,
no second action and persistence $\rho=-1/2$. In the stipulated linear process
the two outcome changes are $(1,-1/2)$. The average change is $1/4$, while
the end-of-period change is $-1/2$. At USD~40 per outcome unit and USD~3
action cost, their net values are respectively USD~7 and USD~$-23$.
Both calculations are correct for their respective targets. An adapter that
returns an average under an endpoint label has reversed the decision without
necessarily changing any estimated coefficient. These values are authored
illustrations, not evidence about actual negative persistence in banking.
\citep{exec16_s2_exact}

The software options address different portions of this problem. Tigramite's
PCMCI-related procedures distinguish temporal conditional-independence,
contemporaneous and latent-variable settings; context and regime extensions
add their own declarations. The inspected JPCMCIplus path treats specified
context classes specially, and the inspected regime method has a restricted
data/mask contract. Those choices cannot be inferred from one generic
``time series'' capability label. CausalAI, VARLiNGAM and DYNOTEARS supply
other comparisons under their stated model and input restrictions. Their
results must be evaluated against a common physical target before speed or
prediction quality is compared. \citep{exec14_pt_tigramite,exec14_pt_causalai,exec14_pt_causalnex,exec13_pydisc_lingam}

Dynamic Bayesian network inference is a useful alternative when the structure
is supplied rather than discovered. The inspected aGrUM unrolling helper and
Mocapy++ inference/parameter-learning paths illustrate that role. A supplied
model can support a transparent finite-state experiment even when no
structure-learning claim is justified. Conversely, a learned temporal graph
need not provide calibrated effects under shared calendar shocks or latent
confounding. The product should expose the supported operation and sampling
boundary rather than merge forecasting, structural discovery and intervention
evaluation into one completion flag. \citep{exec16_ct_agrum,exec16_ct_mocapy,exec16_ct_sdbn}

\section{Interventions provide information only under a mechanism contract}

Regimes are scientifically informative when their relation to the system is
specified. Under a perfect intervention on a target set $I$, the original
mechanisms for variables in $I$ are replaced, while the stipulated mechanisms
outside $I$ remain unchanged. In the inspected interventional-equivalence
formulation, observations are independent across draws but need not be
identically distributed across known targets. A conservative target family
leaves each variable unmanipulated in at least one regime; an observational
regime supplies that property. This condition should not be confused with a
claim that every such experimental design identifies every graph direction.
\citep{exec14_pt_paper_gies}

The likelihood contribution of a variable's original mechanism should not
come from rows on which that mechanism was replaced. This explains why
regime-to-row maps and target indices are part of the score, not incidental
metadata. The inspected pcalg GIES path includes R preprocessing, conversion
to native indices, local-score configuration and a distinction between an
essential graph and a representative. The Python GIES and UTIGSP-related
comparisons likewise need their actual target, test and dependency contracts.
Unknown targets and soft interventions are separate methodological questions;
the assumptions of one known-target procedure cannot be silently transferred
to them. \citep{exec16_ct_pcalg,exec14_pt_gies,exec14_pt_causaldag}

\subsection{Which additional measurements resolve the three-model example}

Return to the three observationally equivalent models $F,R,C$. Replace the
equation for $M$ by the fixed value $m$ and compare $m=1$ with $m=0$.
The population mean changes in coordinate order $(X,M,Y)$ are
\[
 \begin{array}{c|ccc}
  &\Delta X&\Delta M&\Delta Y\\\hline
 F&0&1&1\\
 R&1/2&1&0\\
 C&1/2&1&1
 \end{array}
\]
For $F$, replacing $M$ leaves $X$ unchanged and changes $Y$ by one.
For $R$, it changes $X$ through its half-unit coefficient and leaves the
root $Y$ unchanged. For $C$, both children respond. Thus the joint population
contrast $(\Delta X,\Delta Y)$ distinguishes all three fixed models, while
observing $\Delta Y$ alone distinguishes $R$ from $\{F,C\}$ but does not
separate the two different effects of intervening on $X$.
\citep{exec16_s2_exact}

This is a design insight rather than a finite-sample discovery result. An
actual experiment would observe noisy outcomes, would need sufficient
precision to distinguish relevant contrasts, and would have to justify that
the intervention really replaces only the stipulated mechanism. A marketing
message that also changes expectations through an unmodeled path is not
automatically the intervention in the displayed table. The feasible action,
its target set and the measurements retained must be approved as scientific
design inputs before the comparison can be interpreted that way.

The example also distinguishes effect information from graph orientation
counts. Under the equal model weights and payoffs specified above, current
expected optimal value is USD~2 per decision. Perfect knowledge of the model
would yield $22/3$ dollars, because the offer would be selected only in $F$.
The ideal incremental value is therefore $16/3$ dollars. If only the
population $Y$ response to the intervention on $M$ were revealed, the
$\{F,C\}$ group would have equal conditional weights and expected offer
value $(22-8)/2=7$ dollars. That group has prior probability $2/3$;
the $R$ group would decline. Expected optimal value would be $14/3$,
an improvement of $8/3$ dollars. The omitted $X$ measurement can therefore
matter to the decision even though $Y$ is the final outcome of interest.
\citep{exec16_s2_exact}

These are ideal population-information values under stipulated weights.
They are not trial budgets, package performance or bank ROI. Sampling noise,
intervention cost, delayed outcomes, operational constraints and model
misspecification can lower the value of a proposed experiment. Their role in
the proposal is to identify a concrete question for a cheap generated
prototype: does the complete discovery and reporting path retain the regime
and measurement information needed to make the right conditional decision?

\section{A software portfolio organized by admissible questions}

The detailed survey supports a portfolio rather than a single discovery
winner. For a first causally sufficient independent-observation profile,
pcalg, Tetrad and causal-learn supply transparent constraint-based reference
paths, with score-search candidates kept separate. Their output contracts
should preserve equivalence classes, settings and test failures. For a latent
profile, FCI/RFCI and the inspected native alternatives need their own graph
semantics and computational limits. For temporal observations, a chosen
Tigramite or other temporal path needs a physical clock and explicit sampling
law. For known perfect interventions, a GIES candidate needs accurate targets
and retained unmanipulated mechanisms. These are proposed scoped comparisons,
not permission to enable every operation from a package name.
\citep{exec13_pcalg,exec13_java_tetrad,exec13_pydisc_causallearn,exec13_latentr_primary,exec14_pt_tigramite,exec16_ct_pcalg}

There are also good reasons to choose a different role for a package. A
classification network can be a predictive baseline. A supplied dynamic model
can be an exact or numerical reference. A functional-model optimizer can
generate a candidate under a declared structural family. A commercial editor
can contribute elicitation and model inspection while leaving identification
and calibration to a separate component. The latent commercial comparisons
show why latent-class estimation, parameter learning with missing values and
discovery of an ancestral equivalence class must remain distinct in a
competitive capability matrix. \citep{exec13_java_weka,exec16_ct_agrum,exec16_ct_mocapy,exec13_latent_genie_structure,exec13_latent_mplus}

Direct adoption is most attractive when the existing method and its public
interface already preserve the required scientific object. A language bridge
can retain an established implementation while translating units, names,
endpoint encodings and failures explicitly. A port becomes a separate
procedure to compare, including preprocessing, defaults, stopping and numerical
diagnostics. Source-only reference use is appropriate when rights, dependencies
or essential semantic behavior remain unresolved. Rejecting a proposed use
should specify the mismatch: for example, supplied-model inference does not
satisfy a structure-discovery request. This does not establish that the package
is unsuitable for every other operation. The source-specific reuse sections
develop those decisions beyond this scientific portfolio.

The next inexpensive comparison can use the exact generated mechanisms in
this chapter together with permitted public covariate and calendar structures.
Each engine should receive the same declared question when a matched-method
comparison is intended. A deliberately different method family should instead
be labeled an alternative with its own assumptions. Retain candidate graphs,
equivalence marks, failed requests, conditional effects and decision losses.
Evaluate graph features and the chosen action separately. Include the cost of
scientific preparation and review, rather than reporting only optimizer time.
No public financial table is to be relabeled as intervention truth merely
because its columns look appropriate. \citep{exec16_s2_design}

The wider survey remains bounded by inspected originals and recorded
omissions. Additional versions, less common languages, unknown-target methods,
dependent latent regimes and some commercial capabilities still need
source-specific work. A working numerical example cannot close those search
branches. Equally, an unavailable commercial manual need not block a generated
comparison of the known independent-observation or interventional profiles.
The proposal can advance each admissible question while keeping these limits
visible. Its central product requirement is that the path from observations
through graph uncertainty to the banking decision remain inspectable, with
an explicit explanation when the supported assumptions do not permit an
answer. \citep{exec16_s2_design}
The following implementation studies apply these requirements to concrete
reuse choices. The gCastle and original NOTEARS treatments distinguish
thresholding, numerical stopping and the meaning of a returned matrix.
The Java and R treatments examine what an isolated discovery process would
have to preserve, before the latent-variable comparison changes the graph
class itself. Their source and reference calculations support those choices;
they retain the separate obligation to test an actual installed engine.
\label{exec16:s2-end}
